\section{Implementation Details}
\label{app:implementation}

Experiments were run using a single NVIDIA GeForce GTX 1080 ti, an Intel(R) Core(TM) i7-8700K, on a desktop computer with 16GB of RAM. 
Code is written in python. Packages used include PyTorch \citep{paszke2019pytorch}, scikit-learn \citep{scikit-learn}, Ray \cite{moritz2018ray}, NumPy, SciPy, and Matplotlib. 
We use ray tune \citep{liaw2018tune} with HyperBand Bayesian Optimization \cite{falkner2017combining} search algorithm to optimize our network hyper-parameters. 
The hyper-parameters we consider are accounted for in Table \ref{tab:search_space}. 
The final hyper-parameters used are given in Table \ref{tab:hparams}.
The hyper-parameter optimization objective is the batch-wise Pearson correlation averaged across all outcomes of the validation data for a single dataset realization with random seed 1331.
All experiments reported can be completed in 30 hours using this setup.

\begin{table}[ht]
    \centering
    \begin{tabular}{ll}
        \toprule
        Hyper-parameter & Search Space \\
        \midrule
        hidden units        & tune.qlograndint(32, 512, 32) \\
        network depth       & tune.randint(2, 5) \\
        gmm components      & tune.randint(1, 32) \\
        attention heads     & tune.randint(1, 8) \\
        negative slope      & tune.quniform(0.0, 0.5, 0.01) \\
        dropout rate        & tune.quniform(0.0, 0.5, 0.01) \\
        layer norm          & tune.choice([True, False]) \\
        batch size          & tune.qlograndint(32, 256, 32) \\
        learning rate       & tune.quniform(1e-4, 1e-3, 1e-4) \\
        \bottomrule
    \end{tabular}
    \caption{Hyper-parameter search space}
    \label{tab:search_space}
\end{table}

\begin{table}[ht]
    \centering
    \begin{tabular}{lccc}
        \toprule
        Hyper-parameter & Synthetic & ACCE NN & ACCE Transformer \\
        \midrule
        hidden units        & 96   & 256   & 256 \\
        network depth       & 4     & 3     & 3 \\
        gmm components      & 24    & 24    & 24 \\
        attention heads     & NA    & NA    & 4 \\
        negative slope      & 0.05  & 0.04  & 0.01 \\
        dropout rate        & 0.04  & 0.2  & 0.5 \\
        layer norm          & False  & False  & False \\
        batch size          & 32    & 2048   & 32 \\
        learning rate       & 0.0015  & 1e-4  & 2e-4 \\
        \bottomrule
    \end{tabular}
    \caption{Final hyper-parameters for each dataset/model}
    \label{tab:hparams}
\end{table}

\subsection{Model Architecture}

The general model architecture is shown in \Cref{fig:overcast-architecture}. 
The models are neural-network architectures with two basic components: a feature extractor, $\bm{\phi}(\x; \params)$ ($\bm{\phi}$, for short), and a conditional outcome prediction block $f(\bm{\phi}, \tf; \params)$, or density estimator. 
The covariates $\x$ (represented in \textcolor{blue}{\textbf{blue}}) are given as input to the feature extractor, whose output is concatenated with the treatment $\tf$ (represented in \textcolor{purple}{\textbf{purple}}) and given as input to the density estimator which outputs a Gaussian mixture density $p(\y \mid \tf, \x, \params)$ from which we can sample to obtain samples of the outcomes (represented in \textcolor{red}{\textbf{red}}).
Models are optimized by maximizing the log-likelihood, $\log{p(\y \mid \tf, \x, \params)}$, using mini-batch stochastic gradient descent.

\begin{figure}[ht]
    \centering
    \includegraphics[width=\textwidth]{figures/overcast-architecture}
    \caption{\textbf{Overcast model architecture.} The inputs are represented by circles, in \textcolor{blue}{\textbf{blue}} the covariates, and in \textcolor{purple}{\textbf{purple}} the treatment. In the \textcolor{red}{\textbf{red}} circle is the output of the model, the outcomes distribution. The model has different feature extractors (in \textcolor{green}{green}) for the feed-forward neural network and the transformer. It has a single density estimator (in \textcolor{orange}{orange}).}
    \label{fig:overcast-architecture}
\end{figure}

\subsubsection{Feature extractor}

The feature extractor design is problem and data specific. 
In our case, we look at using both a simple feed-forward neural network and also a transformer. 
The transformer has the advantage of allowing us to model the spatio-temporal correlations between the covariates on a given day using the geographical coordinates of the observations as positional encoding.
This is interesting when studying ACI because confounding may be latent in the relationships between neighboring variables.
Typically, environmental processes (which is one source of confounding) are dependent upon the spatial distribution of clouds, humidity and aerosol, and this feature extractor may capture these confounding effects better. 

\subsection{Density Estimator}

The conditional outcome prediction block, relies on a $n_{\y}$ component Gaussian mixture density represented in \Cref{fig:overcast-gmm}.
It outputs: 
\begin{equation*}
    p(\y \mid \tf, \x, \params) =
    \sum_{j=1}^{n_{\y}} \pitilde_{j}(\bm{\phi}, \tf; \params) \mathcal{N}\left(\y \mid \mut_{j}(\bm{\phi}, \tf; \params), \sigmat^{2}_j(\bm{\phi}, \tf; \params)\right),
\end{equation*}
and 
$$
\mut(\x, \tf; \params) = \sum_{j=1}^{n_{\y}} \pitilde_{j}(\bm{\phi}, \tf; \params)\mut_{j}(\bm{\phi}, \tf; \params),
$$
where $\mathcal{N}(\cdot \mid \mu, \sigma^2)$ denotes a normal distribution with mean $\mu$ and variance $\sigma^2$. 

\begin{figure}
    \centering
    \includegraphics[width=.6\linewidth]{figures/overcast-gmm.pdf}
    \caption{\textbf{Overcast Gaussian mixture model.} The mixing coefficients $\mathbf{\pi}$ are estimated with a linear layer and a SoftMax layer, to obtain $\tilde{\mathbf{\pi}}$, represented in \textcolor{blue}{\textbf{blue}} in the figure. The vector of means of the Gaussian kernels $\tilde{\mathbf{\mu}}$ is obtained by $n_y$ linear layers (in \textcolor{green}{\textbf{green}} in the diagram), whilst the vector of variances $\tilde{\mathbf{\sigma}}$ is obtained by $n_y$ blocks of linear layers and SoftPlus layers (in \textcolor{orange}{\textbf{orange}} in the diagram).}
    \label{fig:overcast-gmm}
\end{figure}